\documentclass[book.tex]{subfiles}

\begin{document}

\chapter{Transformers}

\label{chap:transformer}
\noindent\rule{11cm}{0.4pt} \\
\textbf{Prerequisites:}  \autoref{chap:ml_basics}, \autoref{chap:grad_descent}  \\
\textbf{Difficulty Level:} ** \\
\noindent\rule{11cm}{0.4pt}
\vspace{5mm}

Recurrent neural networks provide powerful tools for working with series, but also have a powerful efficiency disadvantage. Because inputs have to be processed step by step, there are limits to how far RNNs can be parallelized on a modern processor. Transfomers offer an alternative interpretation of sequence handling that allow for more efficient processing.

\begin{figure}
    \centering
    % \includegraphics{figures/Differentiable Models/transformers/transformer_attention_mechanism.jpeg}
    \includegraphics{figures/Differentiable Models/transformers/Transformers-Attention Mechanism.png}
    
    \caption{Transformers use the attention mechanism to estimate the importance of each input in a sequence.}
    \label{fig:transformer}
\end{figure}

\section{Tokenizers and Vocabularies}

Transformers assume that input comes as a sequence $x_1,\dotsc,x_L$. In many applications, we can assume that each input is an integer
\begin{align}
    x_i \in \{1,\dotsc,N\}
\end{align}
where $N$ is a large but finite number. These integers are assumed to index into a vocabulary set $V$, where $i$ represents some element of the vocabulary. Tokenization is the process of transforming an arbitrary input sequence into a sequence of integers drawn from the given vocabulary. An embedding matrix is typically used to associate to each index in the vocabulary an associated high dimensional embedding vector
\begin{align}
    W: \mathbb{R}^{d \times N}
\end{align}

Position encoding is a technique for representing the position $i$ of input $x_i$ in the original sequence to the model. The \lstinline{PositionEncoding} layer provides a way to represent the position of a token by a continuous value. 

\begin{lstlisting}[language=python]
class PositionEncoding(L: $\mathbb{N}$, M: $\mathbb{N}$, $\ell$ : $\mathbb{N}$=1e4):
  def $\lambda$() $\to$ $\mathbb{R}$[L, M]:
    pos = [1..L]
    dim = [1..M]
    phase = pos / ($\ell$ ** (dim/M))   
    where(d % 2 == 0, sin(phase), cos(phase))
\end{lstlisting}


\section{Transformer Equations}

The attention mechanism consists of three terms. First there are a set of query vectors $q_i \in \mathbb{R}^p$ for $i = 1,\dotsc,m$, then a set of key vectors $k_j \in \mathbb{R}^p$ for $j = 1,\dotsc,n$, and then a set of value vectors $v_j \in \mathbb{R}^r$. Here $r$ and $p$ are the sizes of low dimensional learned embeddings. The key $k_j$ and value $v_j$ are conceptually attached to the same input $j$. For query $q_i$ we write the attention mechanism as

\begin{align}
    \textrm{Attn}(q_i, \{k_j\}, \{v_j\}) &= \sum_{j=1}^n \alpha_{ij} v_j \\
    \alpha_{ij} &= \frac{\exp(q_i^Tk_j)}{\sum_{j'=1}^n\exp(q_i^Tk_{j'})}
\end{align}
For self-attention networks, we write
\begin{align}
    q = h_Q(f),\qquad k=h_K(f),\qquad v=h_V(f)
\end{align}
where $h_Q$, $h_K$, and $h_V$ are embedding neural networks.

One of the most powerful properties of the attention mechanism is permutation invariance. Shuffling the order of the inputs does not change the output from the system. Let's look at the code for the attention mechanism.
%\url{https://github.com/fkodom/transformer-from-scratch/blob/main/src/encoder.py}
\begin{lstlisting}[language=python]
class Attention(m : $\mathbb{N}$, n: $\mathbb{N}$, p: $\mathbb{N}$):
  def $\lambda$(query: $\mathbb{R}$[m, p], key: $\mathbb{R}$[n, p], value: $\mathbb{R}$[n, p]): $\mathbb{R}$[m, p]:
    softmax(query*key$^T$/ sqrt(p)) * value

class AttentionHead(D: $\mathbb{N}$, Q: $\mathbb{N}$, K: $\mathbb{N}$):
    q = Linear(D, Q)
    k = Linear(D, K)
    v = Linear(D, K)

    def $\lambda$(query: $\mathbb{R}$[m, D], key: $\mathbb{R}$[n, D], value: $\mathbb{R}$[n, D]):
      Attention(q(query), k(key), v(value))
\end{lstlisting}

In practice, we will often find it useful empirically to run multiple attention heads in parallel, much as we want to learn multiple convolutional filters in parallel.
\begin{lstlisting}[language=python]
class MultiHeadAttention(H: $\mathbb{N}$, D: $\mathbb{N}$, Q: $\mathbb{N}$, K: $\mathbb{N}$):
    heads = [AttentionHead(D, Q, K) for _ in H]
    linear = Linear(H * K, D)

    def $\lambda$(query: $\mathbb{R}$[p], key: $\mathbb{R}$[p, n], value: $\mathbb{R}$[r, d]):
      linear(concat([h(query, key, value) for h in heads]))  
\end{lstlisting}


We introduce a simple feed forward layer for use in more complex architectures.

\begin{lstlisting}[language=python]
def FeedForward(I: $\mathbb{N}$, F: $\mathbb{N}$):
  def $\lambda$(X):
    Linear(I, F) $\circ$ ReLU() $\circ$ Linear(F, I)(X)
\end{lstlisting}
%class Residual(sublayer: Module, M: $\mathbb{N}$, dropout: $\mathbb{R}$):
%    sublayer = sublayer
%    norm = LayerNorm(M)
%    dropout = Dropout(dropout)

%    def $\lambda$(tensors: [Tensor]) $\to$ Tensor:
%      norm(tensors[0] + dropout(sublayer(*tensors)))    

There are several variants of the transformer that have been considered in the literature. In this chapter, we consider the originally proposed sequence-to-sequence encoder-decoder architecture. The encoder encodes the input sequence into a vector representation and the decoder undoes the encoding to make the prediction for the output sequence.

\begin{lstlisting}[language=python]
class EncoderLayer(M: $\mathbb{N}$, H: $\mathbb{N}$, F: $\mathbb{N}$, dropout: $\mathbb{R}$):
    Q = K = max(M // H, 1)
    attention = Residual(
        MultiHeadAttention(H, M, Q, K), M, dropout)
    feed_forward = Residual(
        FeedForward(M, F), M, dropout)

    def $\lambda$(src: $\mathbb{R}$[A, B]):
      src = attention(src, src, src)
      feed_forward(src)

\end{lstlisting}

The \lstinline{Encoder} chains together a sequence of encoding layers.

\begin{lstlisting}[language=python]
class Encoder(L: $\mathbb{N}$, M: $\mathbb{N}$, H: $\mathbb{N}$, F: $\mathbb{N}$, dropout: $\mathbb{R}$):
    layers = [
        EncoderLayer(M, H, F, dropout) for _ in L]

    def $\lambda$(src: $\mathbb{R}$[L, D]):
        src += PositionEncoding(L, D)
        for layer in layers:
          src = layer(src)
        src

\end{lstlisting}

The decoder layer plays a similar role to the encoder layer, but with the critical different that the decoder attempts to match the embedding produced by the encoder to the target sequence.

\begin{lstlisting}[language=python]

        
class DecoderLayer(M: $\mathbb{N}$, H: $\mathbb{N}$, F: $\mathbb{N}$, dropout: $\mathbb{R}$):
    Q = K = max(M // H, 1)
    attention_1 = Residual(
        MultiHeadAttention(H, M, Q, K), M, dropout)
    attention_2 = Residual(
        MultiHeadAttention(H, M, Q, K), M, dropout)
    feed_forward = Residual(
        Feedforward(M, F), M, dropout )

    def $\lambda$(tgt: $\mathbb{R}$[L, D], memory: Tensor) $\to$ Tensor:
        tgt = attention_1(tgt, tgt, tgt)
        tgt = attention_2(tgt, memory, memory)
        feed_forward(tgt)

\end{lstlisting}

The \lstinline{Decoder} chains together a series of decoder layers.

\begin{lstlisting}[language=python]
class Decoder(L: $\mathbb{N}$, M: $\mathbb{N}$, H: $\mathbb{N}$, F: $\mathbb{N}$, dropout: $\mathbb{R}$):
    layers = [
        DecoderLayer(M, H, F, dropout)
        for _ in L]
    linear = Linear(M, M)

    def $\lambda$(tgt: $\mathbb{R}$[L, D], memory: Tensor) $\to$ Tensor:
        tgt += PositionEncoding(L, D)
        for layer in layers:
          tgt = layer(tgt, memory)
        softmax(linear(tgt), dim=-1)

\end{lstlisting}

The \lstinline{Transformer} itself simply chains a decoder and an encoder to prdouce the full sequence-to-sequence architecture.

\begin{lstlisting}[language=python]
class Transformer(NE: $\mathbb{N}$, ND: $\mathbb{N}$, M: $\mathbb{N}$, H: $\mathbb{N}$, F: $\mathbb{N}$, dropout: $\mathbb{R}$):
    encoder = Encoder(NE, M, H, F, dropout)
    decoder = Decoder(ND, M, H, F, dropout)

    def $\lambda$(src: $\mathbb{R}$[L, D], tgt: $\mathbb{R}$[L, D]) $\to$ Tensor:
        decoder(tgt, encoder(src))
\end{lstlisting}

\section{Masked Language Modeling}
Masked language modeling is a type of pretraining in which a random portion of the input data is removed and the model is trained to fill in the missing data. The major advantage of this pretraining is it can be done in a fully self supervised fashion. This and similar pretraining methods have proven highly influential in recent years.
\end{document}